\SourceLeaf{369}%
\noindent die Kosinus der drei Winkel sein, die eine Gerade $0P$ jeweils mit $01,02,03$ bildet. Diese Gerade hat im bewegten Raum nach der Mutation dieselbe Lage wie zuvor. [Wird ferner]
\[
\BareRoot{(aa+\rho\rho)}=k
\]
gesetzt, oder
\[
k=\BareRoot{(aa+bb+cc+dd)},
\]
so ist $k$ die lineare Vergrößerung des mutierten Raums.

Setzt man schließlich
\[
a=k\cos\theta,\qquad \rho=k\sin\theta
\]
so ist $2\theta$ diejenige Drehung um $0P$, durch die in Verbindung mit der Vergrößerung $k$ die Mutation erzeugt werden kann.

\Needspace{12\baselineskip}
\begin{center}\bfseries BEMERKUNG.\end{center}

Die unter [I] zusammengefassten Notizen stehen in einem Handbuch. Die Nummern 1 bis 7 stimmen in Schrift und Tinte genau mit astronomischen Aufzeichnungen aus dem Jahr 1819 überein, die sich auf den vorhergehenden Seiten befinden. Die als Nummer 8 bezeichnete Bemerkung sieht anders aus. Sie trägt inhaltlich den Charakter eines späteren Zusatzes. Diese Annahme wird auch dadurch bestätigt, dass astronomische Notizen, die augenscheinlich gleichzeitig mit ihr niedergeschrieben wurden, aus der Zeit von 1822 bis 1823 stammen. Die Notiz [II] füllt einen einzelnen undatierten Zettel. Nach der Handschrift zu urteilen, stammt sie aus noch späterer Zeit.
\begin{flushright}\textsc{Stäckel.}\end{flushright}
